\section{Buried with Him in Baptism: The First Dying and the Close of Probation}

A Christian does not wait until his final breath to die. In the sixth chapter of Romans, St.~Paul confronts the baptized with a truth so startling that he frames it as an incredulous question:

\begin{quote}
Know you not that all we who are baptized in Christ Jesus are baptized in his death? For we are buried together with him by baptism into death: that, as Christ is risen from the dead by the glory of the Father, so we also may walk in newness of life. For if we have been planted together in the likeness of his death, we shall be also in the likeness of his resurrection. (Rom.\ 6:3--5; cf.\ Col.\ 2:12)
\end{quote}

The baptismal font is a tomb before it is a cradle. Descending into its waters, the old man shaped in Adam is crucified with Christ and buried into His death; rising from the font, the believer already shares the Risen Lord's life hidden with Christ in God (Rom.\ 6:6--11; Col.\ 3:3--4). ``Whether we live, we live unto the Lord: or whether we die, we die unto the Lord. Therefore, whether we live or whether we die, we are the Lord's'' (Rom.\ 14:7--8). Bodily death is thus the Christian's second passage through a door he already entered sacramentally in Baptism.

\subsection{Why the Baptized Still Die}

Yet that doctrine immediately raises an ancient question: if bodily death entered through Adam's sin, and Baptism washes away original sin completely, why are the baptized not made immediately immortal? Why must the saint still grow frail and surrender his breath to the grave?

St.~Augustine answers squarely in Book XIII, chapter 4 of \emph{The City of God}. If bodily immortality followed directly upon leaving the font, ``faith itself would be thereby enervated. For faith is then only faith when it waits in hope for what is not yet seen in substance.'' Crowds would seek Baptism merely for earthly longevity, and the heroic witness of martyrdom would vanish. Instead, divine mercy accomplishes something far greater than erasing the penalty: it converts the ancient punishment of sin into the very armor of holiness:

\begin{quote}
But now, by the greater and more admirable grace of the Saviour, the punishment of sin is turned to the service of righteousness. For then it was proclaimed to man, ``If thou sinnest, thou shalt die;'' now it is said to the martyr, ``Die, that thou sin not.'' Then it was said, ``If ye transgress the commandments, ye shall die;'' now it is said, ``If ye decline death, ye transgress the commandment.'' [\ldots] Thus, by the unutterable mercy of God, even the very punishment of wickedness has become the armour of virtue, and the penalty of the sinner becomes the reward of the righteous.\endnote{St.~Augustine, \emph{The City of God}, XIII.4--5 (Dods, vol.~1), adding: ``God has granted to faith this grace, that death, which is the admitted opposite to life, should become the instrument by which life is reached.''}
\end{quote}

St.~Thomas Aquinas places Augustine's answer inside conformity to Christ. Both sin and mortality are vanquished by the Redeemer, ``but each is done according to the order of Divine wisdom, at a fitting time. Because it is right that we should first of all be conformed to Christ's sufferings, before attaining to the immortality and impassibility of glory, which was begun in Him, and by Him acquired for us.''\endnote{Aquinas, \emph{Summa Theologiae}, I-II, q.~85, a.~5, ad 2; cf.\ III, q.~52, a.~5, ad 2 (citing Rom.\ 8:10).} In St.~Athanasius's phrase, believers ``no longer die the death as before agreeably to the warning of the law; for this condemnation has ceased;'' henceforth we are dissolved according to our mortal nature at the hour God appoints, ``that we may be able to gain a better resurrection.''\endnote{St.~Athanasius, \emph{On the Incarnation}, \S~21.1 (NPNF2 vol.~4, p.~47).}

\subsection{The Single, Irreversible Close of Probation}

If Baptism strips death of hopeless terror, revelation simultaneously surrounds it with awesome moral gravity. Earthly life is the \emph{status viae}---the state of wayfaring and probation---during which the human will chooses freely for or against God under grace. Bodily death is the single, irrevocable boundary where that wayfaring ends.

The Epistle to the Hebrews states the rule with lapidary finality: ``And as it is appointed unto men once to die, and after this the judgment: so also Christ was offered once to exhaust the sins of many'' (Heb.\ 9:27--28; \emph{statutum est hominibus semel mori}). Two words govern the Christian understanding of time: \emph{semel} (``once'') and \emph{post hoc} (``after this''). Each person lives one unrepeatable earthly life and passes once through death. Catholic faith therefore excludes every theory of reincarnation or transmigration of souls, which contradicts both the unique substantial unity of each soul with its own body and the finality of the Paschal mystery.\endnote{See \emph{Catechism of the Catholic Church}, \S~1013 (citing Heb.\ 9:27). St.~Ambrose likewise derides pagan notions that ``transmit the wretched soul into a number of ludicrous monstrosities and figures'' (\emph{On the Decease of His Brother Satyrus}, II.50), and St.~Augustine refutes Platonic cycles of return in \emph{The City of God}, XIII.19.}

Sacred Scripture presses this finality relentlessly. Ecclesiasticus warns that ``in the end of a man is the disclosing of his works'' (Ecclus.\ 11:28--29; cf.\ 7:40). Ecclesiastes sets the immutability of that final orientation before the eye in a stark proverb: ``If the tree fall to the south, or to the north, in what place soever it shall fall, there shall it be'' (Eccles.\ 11:3). Our Lord breaks into the rich fool's plans with the midnight summons---``Thou fool, this night do they require thy soul of thee'' (Luke 12:20; cf.\ 12:16--21)---and commands His disciples to keep watch: ``Wherefore be you also ready, because at what hour you think not the Son of man will come'' (Matt.\ 24:44; cf.\ 24:42; 25:13; Luke 12:35--40).

Why does death fix the will beyond reversal? Neither from arbitrary impatience nor from any defect in divine mercy, but because leaving bodily mutability places the soul in the stable mode of pure spirit. St.~John of Damascus gave the classic formulation:

\begin{quote}
Note, further, that what in the case of man is death is a fall in the case of angels. For after the fall there is no possibility of repentance for them, just as after death there is for men no repentance.\endnote{St.~John of Damascus, \emph{Exposition of the Orthodox Faith} (\emph{De fide orthodoxa}), II.4 (S.~D.~F.~Salmond trans., NPNF, ser.~2, vol.~9, 1898).}
\end{quote}

While living in mortal flesh, the soul moves through shifting passions and remains open to conversion at every breath; once separated from the body, it is fixed in the love or refusal of God in which death found it.

\subsection{The Threshold of Judgment}

Because death closes the time of merit, the departing soul stands at once before Christ in the Particular Judgment and enters immediately into its eternal retribution---direct entrance into heaven, purification in Purgatory for those dying in grace yet still needing cleansing, or immediate descent into hell for those dying in actual mortal sin. Pope Benedict XII's Apostolic Constitution \emph{Benedictus Deus} (1336) and the Council of Florence (\emph{Laetentur caeli}, 1439) solemnly define that immediate retribution; its full unfolding, together with Particular and General Judgment and Purgatory, belongs to Part II of this series.\endnote{Pope Benedict XII, \emph{Benedictus Deus} (29 January 1336), DS 1000--1002; Council of Florence, \emph{Laetentur caeli} (6 July 1439), DS 1304--1306; \emph{Roman Catechism}, pt.~1, Creed Art.~VII (Donovan, pp.~61--63); \emph{Catechism of the Catholic Church}, \S\S~1021--1022; cf.\ Luke 16:22--26; 23:43. Full exposition of Particular and General Judgment and Purgatory belongs to Part II (\emph{The Books Were Opened}), and of Hell and Heaven to Parts III and IV.}

On the nearer side of the threshold, death's finality shatters spiritual lethargy without breeding despair. While breath remains, no sin lies beyond the Blood of Christ: ``Behold, now is the acceptable time: behold, now is the day of salvation'' (2 Cor.\ 6:2). Baptismal union with Christ gives courage before the grave, and the single close of probation lends sobriety to every unrepeatable day of grace.
